\section*{Risk and Environmental Impact Assessment}
\setcounter{section}{6} \setcounter{subsection}{0}
\phantomsection \addcontentsline{toc}{section}{\textit{Risk and Environmental Impact Assessment}}

\noindent
Likelihood $L$, consequence $C$, and risk $R = L \times C$ are scored as in the
QMUL handbook (Sect.~3.6.12). The subsections below give the project-specific
analysis.

\subsection*{1. Risks to successful project completion}

The work combines a custom Gymnasium slicing simulator, Stable-Baselines3 PPO
training, and LangGraph-orchestrated multi-agent LLM policies with external API
calls. The main risks are summarised as Items~R1--R4.

\begin{description}\setlength{\itemsep}{6pt}
\item[R1 --- LLM API outage, quota, or key invalidation]
$L{=}3$, $C{=}3$, $R{=}9$ (significant). Mitigation: use configured retries and
timeouts in \texttt{config.yaml}; fall back to rule-based and PPO baselines;
split LLM batches; document partial results if a full sweep cannot finish.
\item[R2 --- Underestimated experiment and write-up time]
$L{=}3$, $C{=}3$, $R{=}9$ (significant). Mitigation: prioritise core method
comparisons (M1--M5) before optional ablations; reduce exploratory seeds after
pilot runs; freeze scope after an agreed milestone.
\item[R3 --- Implementation or configuration error affecting validity]
$L{=}2$, $C{=}4$, $R{=}8$ (significant). Mitigation: fixed random seeds and
\texttt{config.yaml}; version control; spot-check metrics against baselines;
retain aggregation scripts for reproducibility.
\item[R4 --- Dependency or environment drift]
$L{=}2$, $C{=}2$, $R{=}4$ (moderate). Mitigation: pin dependencies; record Python
and package versions; rerun smoke tests before final data collection.
\end{description}

Malfunctioning software (crashes mid-experiment) is treated as a moderate
operational risk: logging, smaller batches, and resumable scripts where applicable.

\subsection*{2. Risks of harm to people and/or animals}

The project is \emph{computational}: simulation and API-based inference only;
there are no human participants, animals, or field trials. Residual ergonomics
risk from prolonged screen work is rated $L=3$, $C=1$, $R=3$ (low), mitigated by
regular breaks. No further institutional ethics procedures apply beyond normal
safe computing practice.

\subsection*{3. Environmental impact}

\textbf{Energy and emissions.} Training PPO and running experiments uses local
CPU/GPU electricity; LLM calls use remote data-centre energy. We rate avoidable
extra compute and API calls at $L=3$, $C=1$, $R=3$ (low): pilot before large
sweeps, avoid redundant invocations, reuse cached decisions where possible.

\textbf{Waste and hardware.} No chemical or radioactive waste; no dedicated
hardware beyond a standard workstation; end-of-life equipment follows
institutional recycling.

\subsection*{4. Financial loss}

The main exposure is pay-per-use cloud LLM charges at scale. Uncontrolled seeds,
episodes, or debugging runs can inflate cost ($L=4$, $C=2$, $R=8$, significant).
Mitigation: budget cap, monitor token usage, temperature zero to limit retries,
cheaper or offline models for dry-run debugging. No contractual liability to
third parties arises from the simulation-only scope.

\subsection*{Summary}

The highest-priority contingencies are (i) experiments remaining feasible when
the LLM API is constrained, (ii) calendar and API budget control so core
comparisons finish on time, and (iii) validity through configuration control and
sanity checks. Human, animal, and acute environmental harm are negligible for
this software-only study; energy and API cost are managed as low-to-moderate
concerns.
